\subsubsection*{3.1.6. A generalization.}
The preceding discussion generalizes to the case in which
$f:X\to D$ is a morphism such that $0\in X$ is an isolated
non-smooth point \underline{in} $X_0$.  After shrinking $X$ and $D$,
the non-smooth locus of $f$ is then finite over $D$.  The manifold of
vanishing cycles $V$ has a finite set $\Sigma$ of critical points;
$V-\Sigma$ is a differentiable manifold with boundary (its boundary
is compact, although the manifold is not).  The monodromy
transformation $T$ is the identity on the boundary and stabilizes
$\Sigma$.  Finally, Proposition 3.1.5 remains valid: the fibered
homeomorphism $\Psi$ is a diffeomorphism away from the critical locus,
which it carries to the part of $W$ arising from $\Sigma$.

\subsubsection*{3.1.7.}
Proposition 3.1.5 allows us to use the calculations of 1.4.  We obtain:

\begin{enumerate}[label=(\Alph*)]
\item The vanishing-cycle groups $(R^i\Phi(\Lambda))_0$ are the
  reduced cohomology groups, with coefficients in $\Lambda$, of the
  manifold of vanishing cycles $V$.
\item The groups
  $H^i_{\{0\}}(X_0,R\Psi(\Lambda))$ are the cohomology groups with
  proper support of $V^\circ$, with coefficients in $\Lambda$.
\item The variation
  \[
  \operatorname{Var}:\widetilde H^i(V,\Lambda)
  \longrightarrow H_c^i(V^\circ,\Lambda)
  =H^i(V\bmod\partial V,\Lambda)
  \]
  is the classical variation defined by the automorphism $T$ of $V$,
  which is the identity on $\partial V$.
\end{enumerate}

\subsubsection*{3.1.8.}
Assertions (A), (B), and (C) are in fact more elementary than
Proposition 3.1.5.  Here is a direct proof and construction which
spares us from checking that the isomorphisms in (A) and (B) are
independent of the choices made in the construction of 3.1.5.

Let $d\leq r$ and $t\leq\tau(d)$, and let $B$ be the open box
\[
B=\{x\mid\delta(x)<d\ \text{and}\ |f(x)|^2<t\}.
\]
Let $\widetilde D^*$ be a universal covering of $D^*$, and consider
the diagram
\begin{equation}\tag{3.1.8.1}
\begin{tikzcd}[column sep=large,row sep=large]
\{0\} \arrow[r,hook,"\bar\imath"] \arrow[d,equal]
  & \widetilde D_t \arrow[d,"p"]
  & \widetilde D_t^* \arrow[l,hook',"\bar\jmath"'] \arrow[d,"p'"] \\
\{0\} \arrow[r,hook,"i"]
  & D_t
  & D_t^* \arrow[l,hook',"j"']
\end{tikzcd}
\end{equation}
obtained from (1.2.7.2) by restriction to $D_t$.  Let
\begin{equation}\tag{3.1.8.2}
\begin{tikzcd}[column sep=large,row sep=large]
B_0 \arrow[r,hook,"\bar\imath"] \arrow[d,equal]
  & \overline B \arrow[d,"p"]
  & \overline B^* \arrow[l,hook',"\bar\jmath"'] \arrow[d,"p'"] \\
B_0 \arrow[r,hook,"i"]
  & B
  & B^* \arrow[l,hook',"j"']
\end{tikzcd}
\end{equation}
be the inverse image of $f:B\to D$ over (3.1.8.1).  We again write
$f$ for the projection from (3.1.8.2) to (3.1.8.1).

\medskip
\noindent\textup{(A)}
By definition, $R^i\Psi_\eta(\Lambda)_0$ is the direct limit, as $U$
runs through the neighborhoods of $0$ in $X$, of the cohomology groups
\[
H^i(U\times_D\widetilde D^*,\Lambda).
\]
Let $U$ run through the fundamental system of neighborhoods of $0$
formed by boxes $B$ as above.  The fiber product
$U\times_D\widetilde D^*=\overline B^*$ is a bundle with fiber
$V^\circ$ over the contractible space $\widetilde D_t^*$.  Hence the
direct system of the groups $H^i(\overline B^*,\Lambda)$ is constant,
with value
\[
H^i(V^\circ,\Lambda)=H^i(V,\Lambda).
\]
The resulting isomorphisms
\begin{equation}\tag{3.1.8.3}
R^i\Psi_\eta(\Lambda)_0\simeq H^i(V,\Lambda)
\end{equation}
may be refined to the following composite isomorphism (cf. (B)
below):
\begin{equation}\tag{3.1.8.4}
\begin{tikzcd}[column sep=small,row sep=large,cells={nodes={font=\small}}]
R\Psi_\eta(\Lambda)_0
  & R\Gamma(B_0,R\Psi_\eta(\Lambda)) \arrow[l,"\sim"']
  & (Rf_*R\bar\jmath_*\Lambda)_0 \arrow[l,"\sim"']
  & R\Gamma(\widetilde D_t,Rf_*R\bar\jmath_*\Lambda)
      \arrow[l,"\sim"'] \arrow[d,equal] \\
{}
  & R\Gamma(V,\Lambda) \arrow[r,"\sim"]
  & R\Gamma(V^\circ,\Lambda)
  & R\Gamma(\overline B^*,\Lambda) \arrow[l,"\sim"']
\end{tikzcd}
\end{equation}
The complex $R\Phi(\Lambda)_0$ is defined by the distinguished
triangle
\[
\Lambda\longrightarrow(R\Psi_\eta\Lambda)_0
\longrightarrow R\Phi(\Lambda)_0\longrightarrow{}.
\]
Consequently, the groups $R^i\Phi(\Lambda)_0$ are the reduced
cohomology groups of $V$.

\medskip
\noindent\textup{(B)}
Let $\mathfrak F$ be the family of closed subsets of $\overline B^*$
whose closure in $\overline B$ is proper over $\widetilde D$.  The
arrows in the following diagram are isomorphisms:
\begin{equation}\tag{3.1.8.5}
\begin{tikzcd}[column sep=small,row sep=large,cells={nodes={font=\small}}]
R\Gamma_{\{0\}}(R\Psi_\eta(\Lambda))
  \arrow[r,"\sim","\text{\textcircled{1}}"']
  & R\Gamma_c(B_0,R\Psi_\eta(\Lambda))
  & (Rf_!R\bar\jmath_*\Lambda)_0
      \arrow[l,"\sim"',"\text{\textcircled{2}}"] \\
R\Gamma_c(V^\circ,\Lambda)
  & R\Gamma_{\mathfrak F}(\overline B^*,\Lambda)
      \arrow[l,"\sim"',"\text{\textcircled{5}}"]
      \arrow[r,equal,"\text{\textcircled{4}}"']
  & R\Gamma(\widetilde D_t,Rf_!R\bar\jmath_*\Lambda)
      \arrow[u,"\text{\textcircled{3}}","\wr"']
\end{tikzcd}
\end{equation}
Indeed, \textcircled{1} is an isomorphism because $B_0$ is a cone
(3.1.2) and $R^i\Psi_\eta(\Lambda)$ is constant away from $0$;
\textcircled{2} by the proper base-change theorem;
\textcircled{3} because $Rf_!R\bar\jmath_*\Lambda$ is constant away
from $0$; \textcircled{4} is the derived-category version of the
spectral sequence of a composite functor; and \textcircled{5} because
$\widetilde D^*$ is contractible.

The isomorphism (3.1.8.5) supplies the desired isomorphisms
\begin{equation}\tag{3.1.8.6}
R^i\Gamma_{\{0\}}(R\Psi_\eta(\Lambda))
\simeq H_c^i(V^\circ,\Lambda).
\end{equation}
We shall verify (C) in 3.1.10.

